\documentclass[12pt,a4paper]{article}

% --- Packages ---
\usepackage[utf8]{inputenc}
\usepackage[T1]{fontenc}
\usepackage[french]{babel}
\usepackage{geometry}
\geometry{margin=2.5cm}
\usepackage{graphicx}
\usepackage{hyperref}
\hypersetup{
    colorlinks=true,
    linkcolor=blue,
    filecolor=magenta,
    urlcolor=cyan,
    citecolor=blue,
}
\usepackage{enumitem}
\usepackage{booktabs}
\usepackage{xcolor}
\usepackage{titlesec}
\usepackage{fancyhdr}
\usepackage{float}
\usepackage{longtable}
\usepackage{tabularx}
\usepackage{listings}
\usepackage{inconsolata}
\usepackage{tikz}
\usetikzlibrary{arrows.meta, positioning}

\lstset{
    basicstyle=\ttfamily\small,
    breaklines=true,
    frame=single,
    columns=fullflexible,
    keepspaces=true
}

\pagestyle{fancy}
\fancyhf{}
\lhead{Rapport d'Avancement n°3 - SpeechCoach}
\rhead{Master SIE - 2026}
\cfoot{\thepage}

\titleformat{\section}{\large\bfseries\color{blue!70!black}}{}{0em}{}[\titlerule]
\titleformat{\subsection}{\bfseries\color{blue!50!black}}{}{0em}{}
\titleformat{\subsubsection}{\bfseries\color{blue!35!black}}{}{0em}{}

\title{
    \vspace{-2cm}
    \large Master Systèmes Intelligents pour l'Éducation (SIE)\\
    \vspace{0.5cm}
    \huge \textbf{Rapport d'État d'Avancement n°3}\\
    \large \textit{Produit web, architecture technique, RAG/LLM et fine-tuning}
}
\author{\textbf{Étudiant :} Hicham Moussaid \\ \textbf{Encadrant :} Pr. Abdelaaziz Hessane}
\date{Période : après le Rapport n°2 jusqu'à l'état actuel du prototype}

\begin{document}

\maketitle

\section{Synthèse et État Actuel du Projet}
Depuis le dernier rapport, \textbf{SpeechCoach} a évolué d'un moteur analytique vers un \textbf{prototype applicatif complet}. Le projet combine désormais un backend multimodal (FastAPI), un frontend web (Next.js), une base RAG et une reformulation par LLM local.

Le bilan de cette phase est le suivant :
\begin{itemize}
    \item \textbf{Consolidation produit} : authentification, gestion de profil, upload de vidéos et export de rapports PDF/Markdown sont opérationnels.
    \item \textbf{Moteur analytique stable} : les briques déterministes (voix, posture, regard) restent le socle de confiance.
    \item \textbf{Exploration IA avancée} : intégration du RAG pour la pédagogie et expérimentation d'un cycle complet de fine-tuning (QLoRA) pour personnaliser le ton du coaching.
    \item \textbf{Optimisation locale} : adoption d'une architecture hybride (déterministe + Ollama + GGUF) pour garantir la stabilité sur machine personnelle.
\end{itemize}

\section{Rappel du Cadre Initial et Positionnement dans la Roadmap}
Selon la roadmap initiale, les sprints 10 à 12 devaient couvrir :
\begin{itemize}
    \item la construction de l'application ;
    \item la présentation des rapports ;
    \item l'évaluation technique ;
    \item la stabilisation finale.
\end{itemize}

Les travaux réalisés pendant cette période couvrent bien ces objectifs, avec un approfondissement plus important que prévu sur la partie \textbf{coaching IA local}. Là où la roadmap visait initialement un LLM local contrôlé, cette phase a conduit à :
\begin{itemize}
    \item comparer plusieurs familles de modèles ;
    \item revoir plusieurs fois l'architecture de génération ;
    \item tester une approche hybride règles + RAG + LLM ;
    \item lancer une piste de \textbf{fine-tuning supervisé}.
\end{itemize}

\section{Travaux Réalisés Depuis le Deuxième Rapport}

\subsection{Passage à une Architecture Produit Complète}
L'un des changements majeurs est l'évolution du projet d'un pipeline technique vers une architecture applicative complète.

\subsubsection{Backend}
Le backend a été structuré autour de \textbf{FastAPI} \cite{fastapi} pour exposer les fonctionnalités métier, avec une organisation modulaire autour des composants suivants :
\begin{itemize}
    \item authentification et gestion des utilisateurs ;
    \item gestion du profil ;
    \item upload et suivi des vidéos ;
    \item consultation des rapports et de l'historique ;
    \item récupération et mise à jour des feedbacks utilisateurs ;
    \item exposition des fichiers de stockage générés.
\end{itemize}

Le backend s'appuie également sur :
\begin{itemize}
    \item \textbf{SQLModel} pour la persistance ;
    \item \textbf{Celery} \cite{celery} pour l'exécution asynchrone des tâches d'analyse ;
    \item \textbf{Redis} comme broker et backend de tâches ;
    \item un stockage local organisé des fichiers d'entrée et des résultats intermédiaires.
\end{itemize}

\subsubsection{Frontend}
Une interface web complète a été mise en place avec :
\begin{itemize}
    \item \textbf{Next.js} \cite{nextjs} / \textbf{React} \cite{react} ;
    \item \textbf{TypeScript} ;
    \item \textbf{Tailwind CSS} ;
    \item des composants UI modernes (shadcn/base-ui, Framer Motion, Recharts).
\end{itemize}

Cette interface permet aujourd'hui :
\begin{itemize}
    \item l'inscription et la connexion ;
    \item l'import d'une vidéo ou l'enregistrement direct ;
    \item le suivi de traitement ;
    \item la lecture détaillée d'un rapport ;
    \item l'export en PDF et en Markdown ;
    \item la consultation de l'historique et des avis.
\end{itemize}

\subsubsection{Base de données et persistance}
La couche de persistance a été structurée pour supporter un usage produit réel. Le projet repose sur une base relationnelle pilotée par \textbf{SQLModel} \cite{sqlmodel}, avec les principales tables suivantes :

\begin{table}[H]
    \centering
    \small
    \begin{tabularx}{\textwidth}{@{}l >{\raggedright\arraybackslash}X @{}}
        \toprule
        \textbf{Table} & \textbf{Rôle et colonnes principales} \\
        \midrule
        \texttt{user} & compte applicatif ; \texttt{id}, \texttt{email}, \texttt{hashed\_password}, \texttt{is\_active}, \texttt{is\_superuser}, \texttt{is\_verified}. \\
        \texttt{profile} & profil utilisateur ; \texttt{id}, \texttt{user\_id}, \texttt{full\_name}, \texttt{avatar\_url}, \texttt{preferred\_language}, \texttt{experience\_level}, \texttt{current\_goal}, \texttt{weak\_points}. \\
        \texttt{videosession} & session d'analyse ; \texttt{id}, \texttt{user\_id}, \texttt{video\_url}, \texttt{status}, \texttt{duration\_seconds}, \texttt{title}, \texttt{created\_at}. \\
        \texttt{analysisresult} & résultat analytique consolidé ; \texttt{id}, \texttt{video\_session\_id}, \texttt{overall\_score}, \texttt{metrics\_json}. \\
        \texttt{coachingfeedback} & version textuelle du coaching ; \texttt{id}, \texttt{analysis\_result\_id}, \texttt{bilan}, \texttt{conseil}, \texttt{motivation}. \\
        \texttt{platformfeedback} & avis sur l'application ; \texttt{id}, \texttt{user\_id}, \texttt{rating}, \texttt{comments}, \texttt{created\_at}. \\
        \bottomrule
    \end{tabularx}
    \caption{Schéma fonctionnel principal de la base de données.}
\end{table}

Ce schéma distingue clairement :
\begin{itemize}
    \item les données d'identité et de profil ;
    \item les sessions vidéo soumises à l'analyse ;
    \item les résultats analytiques stockés en JSON structuré ;
    \item les sorties textuelles de coaching ;
    \item les retours utilisateur sur l'application.
\end{itemize}

\subsection{Diagramme d'Architecture Simplifié}
Le schéma ci-dessous résume l'architecture technique actuellement retenue.

\begin{figure}[H]
\centering
\begin{tikzpicture}[
    node distance=1.5cm and 1.6cm,
    box/.style={draw, rounded corners, align=center, minimum width=3.4cm, minimum height=1.0cm, fill=blue!5},
    service/.style={draw, rounded corners, align=center, minimum width=3.4cm, minimum height=1.0cm, fill=green!5},
    storage/.style={draw, rounded corners, align=center, minimum width=3.4cm, minimum height=1.0cm, fill=orange!8},
    arrow/.style={-{Latex[length=2.5mm]}, thick}
]

\node[box] (frontend) {Frontend web\\Next.js / React};
\node[box, right=of frontend] (api) {Backend API\\FastAPI};
\node[service, right=of api] (ollama) {LLM local\\Ollama / GGUF};

\node[storage, below=of frontend] (db) {Base de données\\SQLModel / SQLite};
\node[service, below=of api] (celery) {Traitement asynchrone\\Celery + Redis};
\node[service, below=of ollama] (engine) {Moteur analytique\\FFmpeg, Whisper, MediaPipe};

\node[storage, below=of celery] (storage) {Stockage local\\uploads / results};

\draw[arrow] (frontend) -- (api);
\draw[arrow] (api) -- (celery);
\draw[arrow] (celery) -- (engine);
\draw[arrow] (engine) -- (ollama);
\draw[arrow] (api) -- (db);
\draw[arrow] (celery) -- (storage);
\draw[arrow] (engine) -- (storage);
\draw[arrow] (api) -- (ollama);
\end{tikzpicture}
\caption{Architecture simplifiée de SpeechCoach dans son état actuel.}
\end{figure}

\subsection{Fonctionnalités Frontend Implémentées}
L'interface web couvre les principaux cas d'usage via les modules suivants :
\begin{description}
    \item[Authentification :] inscription, connexion JWT et persistance de session.
    \item[Dashboard :] statistiques globales, historique récent, onbroading et progression des scores.
    \item[Studio :] import de fichiers et enregistrement direct (webcam/micro) avec suivi en temps réel.
    \item[Historique :] gestion des sessions (recherche, tri, renommage et suppression).
    \item[Rapport :] verdict visuel, scores par axe, plan d'exercices et export PDF/Markdown.
    \item[Profil \& Paramètres :] gestion de l'identité, avatars, objectifs et sécurité.
\end{description}

\subsection{Stabilisation du Pipeline d'Analyse}
Le moteur multimodal a continué à être consolidé :
\begin{itemize}
    \item audio : transcription, débit, pauses, fillers, énergie ;
    \item vision : présence visage, contact visuel approximatif, mains visibles, intensité gestuelle, qualité visuelle ;
    \item agrégation : scoring par dimension, score global, résumé forces/faiblesses.
\end{itemize}

Les évolutions importantes ont surtout concerné :
\begin{itemize}
    \item le recalibrage des scores ;
    \item la priorisation des recommandations ;
    \item la présentation des résultats de manière plus exploitable pédagogiquement.
\end{itemize}

\subsection{Amélioration du Moteur de Recommandations}
Le moteur de recommandations a été enrichi pour devenir plus explicite et plus priorisé.

Les changements apportés incluent :
\begin{itemize}
    \item ajout d'une structure \texttt{Recommendation} avec catégorie, sévérité, message, \textit{actionable tip} et clé d'exercice ;
    \item amélioration du tri des recommandations selon la faiblesse la plus critique ;
    \item ajustement des poids entre voix, regard, gestuelle et qualité visuelle ;
    \item génération d'un plan de pratique plus cohérent avec les fiches pédagogiques.
\end{itemize}

Un point important de cette période a été la constatation suivante : \textbf{les erreurs utilisateur les plus gênantes n'étaient pas toujours dues au scoring brut, mais souvent à l'ordre de priorité présenté dans le rapport}. Une part significative du travail a donc porté sur le recalibrage de cette hiérarchie.

\section{Stack Technique Actuelle}
\begin{table}[H]
    \centering
    \begin{tabular}{@{}ll@{}}
        \toprule
        \textbf{Couche} & \textbf{Technologies principales} \\
        \midrule
        Ingestion vidéo & FFmpeg \cite{ffmpeg}, OpenCV \cite{opencv_library} \\
        ASR / audio & Faster-Whisper \cite{radford2022robust,ctranslate2}, Librosa \cite{mcfee2015librosa} \\
        Vision & MediaPipe \cite{mediapipe2023} \\
        Recherche sémantique & Sentence-Transformers \cite{reimers-2019-sentence-bert}, FAISS \cite{johnson2019billion} \\
        Backend API & FastAPI \cite{fastapi}, SQLModel, Celery \cite{celery}, Redis \\
        Frontend & Next.js \cite{nextjs}, React \cite{react}, TypeScript, Tailwind CSS \\
        LLM local (tests) & Ollama \cite{ollama}, Qwen, Llama, Mistral \\
        Fine-tuning & Transformers \cite{transformers}, PEFT \cite{peft}, TRL, QLoRA \cite{qlora} \\
        \bottomrule
    \end{tabular}
    \caption{Stack technique consolidée dans l'état actuel du prototype.}
\end{table}

\section{Travaux sur la Restitution du Rapport}
L'un des objectifs majeurs de cette phase a été d'améliorer la \textbf{lisibilité pédagogique} du rapport.

Auparavant, le système exposait trop directement les métriques techniques. Le rapport a été réorganisé selon une logique plus lisible :
\begin{itemize}
    \item verdict global ;
    \item bilan ;
    \item priorité ;
    \item encouragement ;
    \item scores par axe ;
    \item points forts ;
    \item axes de progression ;
    \item détails techniques dans une section repliable.
\end{itemize}

Cette refonte a permis de mieux séparer :
\begin{itemize}
    \item la lecture rapide pour l'utilisateur ;
    \item l'explication technique détaillée pour l'analyse.
\end{itemize}

\section{Travaux sur RAG et Connaissances Pédagogiques}
La base de connaissances a été structurée sous forme de fiches d'exercices portant notamment sur :
\begin{itemize}
    \item débit trop rapide ;
    \item débit trop lent ;
    \item fillers et hésitations ;
    \item contact visuel ;
    \item posture figée ;
    \item gestes trop agités ;
    \item cadrage ;
    \item luminosité et netteté.
\end{itemize}

Le RAG a été exploité principalement pour :
Le RAG a été exploité principalement pour :
\begin{itemize}
    \item récupérer les fiches pédagogiques les plus pertinentes ;
    \item enrichir le contexte pédagogique transmis au module de coaching ;
    \item fournir une base explicative complémentaire au rapport.
\end{itemize}
\subsection{Constat important sur le RAG}
Cette période a permis de clarifier un point essentiel : le RAG est \textbf{utile}, mais pas partout.

Il s'est révélé particulièrement utile pour :
\begin{itemize}
    \item les exercices conseillés ;
    \item les explications pédagogiques ;
    \item l'enrichissement du plan de pratique.
\end{itemize}

En revanche, son usage comme source principale pour générer directement le bilan, la priorité et l'encouragement s'est montré plus fragile, car les passages récupérés pouvaient polluer la formulation finale et faire dériver le ton du rapport. Dans l'architecture actuelle, ces trois champs restent pilotés par le moteur déterministe et reformulés par le LLM, avec un contexte RAG seulement en appui.

\section{Expérimentations LLM : Modèles Essayés}
Une grande partie du travail récent a consisté à expérimenter plusieurs modèles locaux pour générer les champs courts du rapport.

\subsection{Modèles testés}
Les modèles suivants ont été testés à différentes étapes :
\begin{itemize}
    \item \textbf{Llama 3.2 1B / 3B} ;
    \item \textbf{Llama 3.1 8B} ;
    \item \textbf{Qwen 2.5 3B Instruct} ;
    \item \textbf{Qwen 2.5 7B Instruct} ;
    \item \textbf{Mistral 7B / latest}.
\end{itemize}

\subsection{Constats empiriques}
Les observations principales ont été les suivantes :
\begin{itemize}
    \item les petits modèles sont rapides mais fragiles sur la cohérence logique ;
    \item les modèles plus gros ne garantissent pas une meilleure qualité en français contraint ;
    \item \textbf{Qwen 2.5 3B Instruct} s'est révélé le meilleur compromis parmi les modèles testés pour ce besoin précis ;
    \item \textbf{Mistral} a montré un comportement technique correct mais une qualité linguistique insuffisante pour des champs courts en français ;
    \item \textbf{Llama 3.x} n'a pas produit des résultats suffisamment propres et stables sur les sorties très contraintes.
\end{itemize}

\section{Problématique Centrale Rencontrée avec les LLM}
La difficulté majeure n'a pas été la génération longue, mais au contraire la production de \textbf{trois textes très courts et très contraints} :
\begin{itemize}
    \item \texttt{bilan\_global} ;
    \item \texttt{point\_prioritaire} ;
    \item \texttt{encouragement}.
\end{itemize}

Ce type de génération s'est révélé plus difficile que prévu pour un LLM local, car il faut simultanément :
\begin{itemize}
    \item respecter strictement les métriques ;
    \item garder un ton simple et professionnel ;
    \item éviter le langage générique ;
    \item rester court ;
    \item ne pas halluciner ;
    \item ne pas changer de registre (\og tu \fg{} / \og vous \fg{}).
\end{itemize}

\subsection{Types d'erreurs réellement observées}
Les échecs observés au fil des tests incluent :
\begin{itemize}
    \item fuite de prompt ;
    \item mélange français / anglais ;
    \item encouragements génériques ou creux ;
    \item tournures traduites ou peu naturelles ;
    \item usage accidentel du tutoiement ;
    \item conseils trop verbeux ou mal reliés à l'axe prioritaire ;
    \item contamination par la fiche RAG ;
    \item JSON invalide ou clés incorrectes.
\end{itemize}

\section{Solutions Essayées pour Corriger les Sorties LLM}

\subsection{Prompt Engineering}
Un travail important d'itération a été mené sur les prompts :
\begin{itemize}
    \item séparation des champs ;
    \item tentative de sortie JSON unique ;
    \item consignes de style plus strictes ;
    \item exemples positifs et négatifs ;
    \item limitation explicite du nombre de phrases ;
    \item interdiction des formulations lyriques ou génériques ;
    \item consignes explicites de vouvoiement ;
    \item baisse de la température ;
    \item validation du JSON en sortie.
\end{itemize}

\subsection{Validation et garde-fous}
Des mécanismes de sécurité ont été ajoutés :
\begin{itemize}
    \item nettoyage agressif des sorties ;
    \item détection de tutoiement ;
    \item détection de formulations interdites ;
    \item fallback si la phrase générée est trop faible ;
    \item validation du schéma attendu ;
    \item contrôle du nombre de phrases.
\end{itemize}

\subsection{Recentrage architectural}
Un résultat important de cette phase a été la conclusion suivante :
\begin{quote}
Le LLM est utile comme \textbf{couche d'enrichissement ou de reformulation}, mais il n'est pas suffisamment fiable pour être la source unique de vérité du rapport.
\end{quote}

Cela a conduit à une architecture de plus en plus hybride :
\begin{itemize}
    \item déterministe pour le scoring, les forces, les faiblesses et les recommandations ;
    \item base de connaissances et mapping déterministe pour les exercices, avec récupération sémantique pour enrichir certaines explications ;
    \item LLM comme reformulation ou enrichissement contrôlé.
\end{itemize}

\section{Travaux de Fine-Tuning}

\subsection{Motivation}
Face aux limites du prompting seul, une nouvelle piste a été explorée : le \textbf{fine-tuning supervisé} d'un petit modèle local pour l'adapter plus précisément au style de SpeechCoach.

\subsection{Objectif du fine-tuning}
L'objectif n'était pas d'améliorer le scoring ni le raisonnement métier, mais d'apprendre au modèle à produire des sorties plus stables pour :
\begin{itemize}
    \item le respect du schéma JSON ;
    \item la qualité du français ;
    \item la cohérence de ton ;
    \item le respect de la structure des trois champs.
\end{itemize}

\subsection{Méthode de constitution du dataset}
Le dataset n'a pas été construit comme un simple export brut des rapports. Une phase de structuration a été nécessaire.

\subsubsection{Source des données}
Les exemples ont été construits à partir de rapports réels générés par le projet, en exploitant notamment :
\begin{itemize}
    \item les scores ;
    \item les forces et faiblesses ;
    \item les recommandations ;
    \item les \textit{actionable tips}.
\end{itemize}

\subsubsection{Format et structure des données}
Le dataset a d'abord été construit sous un format simple \textbf{input/target} (score, forces, faiblesses $\rightarrow$ bilan, priorité, encouragement) pour faciliter la validation humaine, avant d'être converti au format final \textbf{messages} (type chat) pour l'entraînement.

\subsubsection{Format de type chat (Entraînement)}
Le format final utilisé pour le fine-tuning (compatible avec les modèles instruct et les tokenizers modernes) est le suivant :

\begin{lstlisting}[language={}]
{
  "messages": [
    {
      "role": "system",
      "content": "Tu es un coach de prise de parole. Tu dois repondre uniquement en JSON valide..."
    },
    {
      "role": "user",
      "content": "Score global: 69\nVoix: 6.1/10\n..."
    },
    {
      "role": "assistant",
      "content": "{\"bilan_global\":\"...\",\"point_prioritaire\":\"...\",\"encouragement\":\"...\"}"
    }
  ]
}
\end{lstlisting}

\subsubsection{Stratégie de génération des données}
Deux méthodes ont été utilisées :
\begin{itemize}
    \item extraction de cas cohérents depuis les rapports existants ;
    \item création de cas synthétiques mais alignés avec la logique du moteur afin de couvrir davantage de scénarios.
\end{itemize}

\subsubsection{Scénarios couverts}
Le dataset a été enrichi pour inclure plusieurs familles de situations :
\begin{itemize}
    \item débit trop rapide ;
    \item débit trop lent ;
    \item fillers et hésitations ;
    \item regard caméra faible ou moyen ;
    \item gestuelle trop figée ;
    \item gestes trop agités ;
    \item cadrage instable ;
    \item image floue ;
    \item éclairage faible.
\end{itemize}

\subsection{Nettoyage et itérations du dataset}
Le travail sur le dataset a demandé plusieurs itérations :
\begin{itemize}
    \item première version manuelle de quelques exemples ;
    \item extension progressive à 20 puis 50 exemples ;
    \item tentative d'extension automatique vers 300 exemples via IA ;
    \item nettoyage important pour supprimer les doublons, les catégories non alignées avec le moteur et les formulations artificielles ;
    \item version consolidée d'environ \textbf{159 exemples exploitables}.
\end{itemize}

\subsection{Outils et méthode de fine-tuning}
Le fine-tuning a été réalisé sur \textbf{Google Colab} avec GPU gratuit, en utilisant :
\begin{itemize}
    \item \textbf{Transformers} \cite{transformers} ;
    \item \textbf{PEFT} \cite{peft} ;
    \item \textbf{TRL} ;
    \item \textbf{QLoRA} \cite{qlora}.
\end{itemize}

Le modèle ciblé pour cette première expérience a été :
\begin{itemize}
    \item \textbf{Qwen/Qwen2.5-3B-Instruct} \cite{qwen25}.
\end{itemize}

\subsection{Résultats obtenus}
Les résultats ont montré des progrès réels :
\begin{itemize}
    \item meilleure stabilité du schéma JSON ;
    \item bonne génération des bonnes clés ;
    \item meilleure cohérence entre entrée et sortie ;
    \item réduction des fuites de prompt ;
    \item amélioration du respect de la structure globale.
\end{itemize}

Cependant, des limites sont restées présentes :
\begin{itemize}
    \item certains champs ne respectaient pas toujours exactement le nombre de phrases attendu ;
    \item l'encouragement restait le champ le plus fragile ;
    \item certains cas rares produisaient encore des formulations imparfaites.
\end{itemize}

\section{Difficultés de Déploiement du Modèle Fine-Tuné}
Le fine-tuning a donné de bons résultats en environnement Colab, mais son \textbf{déploiement local} a révélé de nouvelles contraintes.

\subsection{Problèmes rencontrés}
\begin{itemize}
    \item nécessité de charger à la fois le \textbf{modèle de base} et l'\textbf{adapter LoRA} ;
    \item poids important de la base Qwen 2.5 3B ;
    \item temps de téléchargement long ;
    \item consommation disque importante ;
    \item chargement CPU lourd sur machine locale ;
    \item temps d'enrichissement fortement allongé ;
    \item pics de mémoire vive dans le worker Celery ;
    \item exécution instable en environnement local Windows/CPU.
\end{itemize}

\subsection{Pourquoi l'approche \textit{base model + adapter} a montré ses limites}
L'approche initiale consistait à charger le modèle de base via Hugging Face, puis à y rattacher l'adapter fine-tuné avec \textbf{PEFT}. Cette logique est correcte du point de vue théorique, mais elle s'est révélée difficile à soutenir dans le contexte applicatif réel du projet.

Le backend exécute l'enrichissement dans un worker Celery local, sur CPU, en parallèle du pipeline d'analyse. Pendant les essais, le chargement du modèle de base et de l'adapter à l'intérieur du même processus a produit :
\begin{itemize}
    \item des temps de chargement élevés ;
    \item une forte occupation mémoire ;
    \item des risques de plantage du worker ;
    \item une très mauvaise expérience de déploiement local.
\end{itemize}

\subsection{Solutions étudiées face à ce problème}
Plusieurs pistes ont été explorées :
\begin{itemize}
    \item conserver le modèle fine-tuné tel quel et compenser ses défauts par du fallback ;
    \item tester un modèle plus petit, notamment \textbf{Qwen 2.5 1.5B Instruct} ;
    \item pré-télécharger localement le modèle de base ;
    \item déplacer la charge d'inférence en dehors du worker Python principal ;
    \item fusionner le modèle de base et l'adapter dans un seul artefact prêt à servir.
\end{itemize}

Les essais sur \textbf{Qwen 2.5 1.5B Instruct} ont été utiles d'un point de vue expérimental, mais ils ont également montré une baisse sensible de la qualité sémantique et linguistique. Cette piste est donc restée intéressante, sans devenir le choix principal pour le produit.

\subsection{Fusion des poids : pourquoi et comment}
La solution retenue a consisté à \textbf{fusionner} les poids appris par l'adapter LoRA avec un modèle de base compatible afin d'obtenir un \textbf{modèle unique}.

La combinaison se fait en trois étapes :
\begin{enumerate}
    \item on part d'un modèle de base ;
    \item on y applique les poids appris par l'adapter LoRA ;
    \item on exporte ensuite un modèle fusionné unique.
\end{enumerate}

L'intérêt de cette fusion est double :
\begin{itemize}
    \item supprimer le chargement séparé du modèle de base et de l'adapter au runtime ;
    \item transformer le résultat du fine-tuning en un artefact unique, plus simple à déployer.
\end{itemize}

\subsection{Quantification et passage au format GGUF}
Après fusion, le modèle a été converti dans un format plus adapté au déploiement local :
\begin{itemize}
    \item format \textbf{GGUF} ;
    \item quantification 4 bits de type \texttt{Q4\_K\_M}.
\end{itemize}

Le but de cette étape était de réduire :
\begin{itemize}
    \item la taille disque ;
    \item la pression mémoire ;
    \item le coût d'inférence sur machine locale CPU.
\end{itemize}

\subsection{Pourquoi Ollama a été retenu pour le runtime}
Une fois le modèle fusionné et quantifié, il restait à choisir une stratégie d'inférence stable. Le projet a finalement retenu \textbf{Ollama} \cite{ollama}.

Ce choix s'explique par plusieurs raisons pratiques :
\begin{itemize}
    \item Ollama encapsule l'inférence locale et évite de charger de gros objets PyTorch directement dans le worker applicatif ;
    \item l'isolation mémoire est meilleure que dans l'approche initiale ;
    \item l'intégration côté backend devient simple, sous forme d'appel local au service Ollama ;
    \item l'ensemble est beaucoup plus robuste en usage produit local.
\end{itemize}

\subsection{Architecture LLM actuellement retenue}
L'architecture actuelle du coaching peut être résumée ainsi :
\begin{itemize}
    \item scoring, hiérarchisation des priorités et recommandations : \textbf{déterministes} ;
    \item exercices et plan de pratique : \textbf{base de connaissances + mapping déterministe par \texttt{exercise\_key}} ;
    \item récupération de fiches pédagogiques complémentaires : \textbf{RAG} ;
    \item reformulation des champs courts (\texttt{bilan\_global}, \texttt{point\_prioritaire}, \texttt{encouragement}) : \textbf{modèle local fusionné, quantifié et servi par Ollama, avec contexte RAG en appui} ;
    \item validation, nettoyage et fallback : \textbf{déterministes}.

\section{Tableau de Synthèse : Implémenté / Restant à Faire}
Afin d'éviter de disperser l'état d'avancement dans plusieurs sections, le tableau suivant résume l'état actuel des principales briques du projet.

\begin{table}[H]
    \centering
    \begin{tabular}{@{}p{5.2cm}p{3.2cm}p{6cm}@{}}
        \toprule
        \textbf{Élément} & \textbf{Statut} & \textbf{Commentaire} \\
        \midrule
        Authentification (inscription / connexion) & Implémenté & Parcours frontend disponible avec session JWT. \\
        Profil utilisateur & Implémenté & Modification du profil, avatar, langue, niveau et objectif. \\
        Changement de mot de passe & Implémenté & Disponible dans les paramètres. \\
        Suppression du compte & Partiel / à finaliser & Présentée côté interface, non activée côté backend. \\
        Studio import vidéo & Implémenté & Téléversement, progression, lancement de l'analyse. \\
        Studio enregistrement direct & Implémenté & Webcam + microphone + prévisualisation. \\
        Historique des analyses & Implémenté & Recherche, tri, renommage, suppression. \\
        Lecture du rapport web & Implémenté & Verdict, scores, priorités, exercices, transcription. \\
        Export PDF / Markdown & Implémenté & Deux formats de restitution disponibles. \\
        Scoring déterministe & Implémenté & Partie la plus stable du système. \\
        RAG d'exercices & Implémenté & Utile pour enrichir les plans de pratique. \\
        LLM local non fine-tuné & Expérimenté & Qualité variable sur les sorties courtes. \\
        Fine-tuning QLoRA & Expérimenté & Résultats encourageants, mais déploiement initial trop lourd. \\
        Runtime GGUF + Ollama & Implémenté & Option actuellement la plus réaliste localement. \\
        Évaluation utilisateur formelle & À approfondir & Encore à structurer pour la phase mémoire finale. \\
        \bottomrule
    \end{tabular}
    \caption{Vue synthétique de l'état d'avancement fonctionnel et technique.}
\end{table}

\section{Bilan, Enseignements et Perspectives}
L'état actuel du projet et les enseignements tirés de cette phase peuvent être synthétisés ainsi :

\begin{itemize}
    \item \textbf{Solidité du socle :} le pipeline multimodal et le scoring déterministe sont stables et constituent la source de vérité du système.
    \item \textbf{Apports du RAG :} la base de connaissances pédagogiques enrichit efficacement les plans de pratique sans complexité excessive.
    \item \textbf{Leçons sur le LLM :} la reformulation par LLM local (Ollama/GGUF) est utile mais doit être strictement encadrée par des garde-fous déterministes pour éviter les hallucinations et les problèmes de ton. 
    \item \textbf{Fine-tuning :} le cycle QLoRA a prouvé son utilité pour le respect des formats (JSON) et du style, bien que son déploiement nécessite une fusion de poids et une quantification GGUF pour rester réaliste sur CPU.
    \item \textbf{Problèmes résolus :} la friction entre qualité de génération et temps d'exécution a été traitée par une architecture hybride (logique métier + LLM asynchrone).
\end{itemize}

Les perspectives pour les prochains sprints concernent désormais :
\begin{itemize}
    \item \textbf{Adaptabilité multi-appareils :} le système doit encore être recalibré pour des enregistrements issus de smartphones, tablettes et PC, car des signaux comme le regard caméra, le pitch et le yaw ne se lisent pas de la même manière selon la position de l'appareil.
    \item \textbf{Plan de progression :} le format actuel sur 7 jours fonctionne comme première version, mais reste trop statique. Il faudra le faire évoluer vers un plan plus court, plus adaptatif et plus cohérent avec le profil et le niveau de l'utilisateur.
    \item \textbf{Responsivité et interface :} l'interface devra être davantage adaptée aux différents écrans, avec un travail complémentaire sur la qualité visuelle, l'ergonomie et l'ajout de nouveaux graphiques.
    \item \textbf{Fonctionnalités produit restantes :} certaines briques doivent encore être ajoutées ou finalisées, notamment la vérification d'e-mail et quelques éléments de robustesse dans le parcours utilisateur.
    \item \textbf{Évaluation :} une étude utilisateur plus formelle restera nécessaire pour la phase mémoire finale.
\end{itemize}

\section{Conclusion}
Cette phase du projet a été particulièrement riche, car elle a permis de dépasser une vision simpliste du type \og ajouter un LLM = améliorer automatiquement le coaching \fg. L'expérience montre au contraire que :
\begin{itemize}
    \item les composantes déterministes bien conçues restent essentielles ;
    \item le RAG apporte une valeur réelle mais ciblée ;
    \item la qualité d'un LLM local dépend fortement du contexte, de la langue, du format attendu et des contraintes de déploiement ;
    \item le fine-tuning est une piste prometteuse, mais qui doit être évaluée non seulement en qualité, mais aussi en coût d'intégration.
\end{itemize}

Le projet SpeechCoach dispose désormais d'une base technique et produit beaucoup plus mature qu'au précédent rapport. Les prochaines étapes porteront donc moins sur l'ajout de nouvelles briques que sur l'arbitrage entre \textbf{robustesse, qualité du coaching généré et coût d'exécution}.

\section{Références Bibliographiques}
\nocite{*}
\bibliographystyle{ieeetr}
\bibliography{references}

\end{document}

